% PlanetMath proposal draft for arXiv:2602.09371
% Source corpora: arXiv math.CT eprints, storage/mo-processed-gpu, storage/math-processed-gpu
\section*{Acyclic complexes of FP-injectives over Ding-Chen rings}
\subsection*{Source}
\begin{itemize}
  \item arXiv:2602.09371, James Gillespie, ``Acyclic complexes of FP-injective modules over Ding-Chen rings'' (2026-02-15)
  \item Cross-reference MO 325938 on fp-injective (absolutely pure) modules and concrete examples
  \item Cross-reference Math.SE 3126464 on when FP-injective, pure-injective modules are actually injective
\end{itemize}
\subsection*{Synopsis}
Gillespie introduces a method for merging cotorsion pairs to build abelian model structures and applies it to left coherent (in particular Ding-Chen) rings. The resulting model structure has fibrant objects generated by weakly Ding injective modules, characterized as cycles of specific acyclic complexes of FP-injective modules. Over Ding-Chen rings these coincide with cycles of all acyclic FP-injective complexes, leading to a new description of the stable module category $\mathrm{Stmod}(R)$ via \emph{Gorenstein FP-pro-injective} modules—cycles of totally acyclic complexes of FP-projective-injective modules. A separate application shows that any complete cotorsion pair (even non-hereditary) lifts to an abelian model for the derived category. This gives practical tools for studying fp-injective phenomena discussed in MO 325938 and Math.SE 3126464.
\subsection*{MathOverflow cues}
\begin{description}
  \item[MO 325938] asks for definitions and examples of fp-injective modules; Gillespie’s characterization via acyclic complexes provides explicit constructions and relates them to Ding injective modules.
\end{description}
\subsection*{Math.SE anchors}
\begin{description}
  \item[Math.SE 3126464] inquires when FP-injective plus pure-injective implies injective; the paper’s analysis of weakly Ding injectives clarifies such implications under coherence assumptions.
\end{description}
\subsection*{Encyclopedia outline}
\begin{enumerate}
  \item \textbf{Background on FP-injectives and Ding-Chen rings.} Review definitions, examples (linking to MO 325938), and relevance of coherence.
  \item \textbf{Merging cotorsion pairs.} Explain the new method to obtain abelian model structures and specify the key cotorsion pairs involved.
  \item \textbf{Weakly Ding injectives and fibrant objects.} Present the equivalences between weakly Ding injective modules, cycles of certain FP-injective complexes, and conditions tied to Math.SE 3126464.
  \item \textbf{Gorenstein FP-pro-injective modules.} Define these modules, describe their characterization via totally acyclic FP-projective-injective complexes, and explain the resulting model for $\mathrm{Stmod}(R)$.
  \item \textbf{Applications to derived categories.} Outline how non-hereditary cotorsion pairs still yield abelian models for $\mathbf{D}(R)$.
\end{enumerate}
\subsection*{Action items}
\begin{itemize}
  \item Summarize the construction of the merged model structure with clear diagrams of the cotorsion pairs.
  \item Work out an explicit example over a Ding-Chen ring (e.g., an $n$-FC ring) showing the new fibrant objects.
  \item Detail the definition and properties of Gorenstein FP-pro-injective modules and relate them to classical Gorenstein injectives.
  \item Provide references and cross-links for FP-injective basics, Ding injective modules, and stable module categories.
\end{itemize}
